\section*{Exercises}
\begin{exercise}\label{ex:18.1}Starting from the one-step estimate $e_{n+1}\le qe_n+\eps_n$, derive bounds for $e_1$, $e_2$ and $e_3$ in terms of $e_0$ and the update errors, without using summation notation. Then, in the bound for $e_n$ in part~1 of Theorem~\ref{thm:inexact}, find the power of $q$ that multiplies the earliest update error $\eps_0$ and the power that multiplies the latest one, $\eps_{n-1}$, and explain why these are the right powers.\end{exercise}
\begin{exercise}\label{ex:18.2}Let $0<q<1$ and $\eta\ge0$, and define $y_0=0$ and $y_{n+1}=qy_n+\eta$ for every $n\in\NN$. Prove by induction that $y_n=\eta(1-q^n)/(1-q)$ for every $n\in\NN$. Then take $\eta>0$, and explain why the factor $1/(1-q)$ in part~2 of Theorem~\ref{thm:inexact} cannot be replaced by any smaller number $K$.\end{exercise}
\begin{exercise}\label{ex:18.3}Let $\eta>0$, let $T(x)=0$ for every $x\in\RR$, and let $y_n=(-1)^n\eta$ for every $n\in\NN$. Check that all the hypotheses of Theorem~\ref{thm:inexact} hold with update errors $\eps_n=\eta$, and check directly that the bound of part~2 holds. Then prove that the sequence $y_0,y_1,y_2,\ldots$ does not converge, and explain why a valid error bound does not settle whether a sequence converges.\end{exercise}
\begin{exercise}\label{ex:18.4}Suppose that the update errors in Theorem~\ref{thm:inexact} satisfy $\eps_j\le aq^j$ for every $j\in\NN$, where $a\ge0$ and $q$ is the contraction factor. Derive the bound $e_n\le q^n e_0+anq^{n-1}$ for every $n\ge1$, and prove that this upper bound tends to zero as $n$ grows. (Hint: the factor $n$ grows, so it is not enough that $q^{n-1}\to0$. Put $r=(1+q)/2$ and $t=q/r$, and compare $nt^{n-1}$ with $1+t+\cdots+t^{n-1}$.)\end{exercise}
\begin{exercise}\label{ex:18.5}A finite set $L$ of students and a finite set $R$ of tasks are given, and each student accepts some of the tasks, as in Chapter~\ref{ch:hall}. For a set $S\subseteq L$, $N(S)$ is the set of tasks accepted by at least one student in $S$. Each task $r\in R$ now has a \emph{capacity} $c_r$, which is a nonnegative integer. An \emph{assignment with capacities} gives every student one task that the student accepts, in such a way that each task $r$ is given to at most $c_r$ students. Prove that an assignment with capacities exists if and only if
\[
 |S|\le\sum_{r\in N(S)}c_r\qquad\text{for every set }S\subseteq L.
\]
When every capacity equals $1$, this is Hall's theorem.\end{exercise}
\begin{exercise}\label{ex:18.6}Write a short record of the investigation of the biased iteration $y_{n+1}=(3y_n+1)/4+1/100$, $y_0=0$, of Section~\ref{ch18:sec:experiment}. It should contain the exact question, the background the argument needs, an outline of the proof, the counterexample that refutes the first guess, one numerical check done in exact arithmetic, and an honest statement of what is new. Then state one further question whose answer does not follow from the proofs in this chapter.\end{exercise}
